\section{Microphysics}\label{sec:microphys}
{\bf \Large
\begin{tabular}{ccc}
\hline
  Corresponding author & : & Yuta Kawai\\
\hline
\end{tabular}
}

To represent cloud-microphysics processes, 
we provide two microphysics schemes using SCALE library: 
three-class one-moment bulk and six-class one-moment bulk schemes.

In addition, 
a large-scale condensation scheme can be used. 
It is useful to investigate a feedback mechanism between the equations of motion and moisture within simplified physics.

\subsection{Three-class one-moment bulk scheme}
%%
As a three-class one-moment bulk microphysics scheme, 
the Kessler parameterization \citep{kessler_1969} can be used. 
In this scheme, 
the distribution function of particle size is expressed only by mass concentration. 
Considering two categories of water in cloud and rain, 
the ratios of the densities of cloud and rain to total air density are calculated.

The details of formulation are found in Section 8.3.1 of SCALE description \citep{scale_library}.

\subsection{Six-class one-moment bulk scheme}
%%
A six-class one-moment bulk microphysics scheme is based on \citet{tomita_2008}. 
This scheme calculates mass exchange among six categories of water substances, 
which are water vapor, cloud water, rain, cloud ice, snow, and graupel.  
Although the scheme largely follows \citet{lin_etal_1983}, 
the original method of \citet{lin_etal_1983} is modified as follows: 
Both cloud water and cloud ice are generated only by a saturation adjustment process, 
and wet growth process of graupel is omitted.
\citet{tomita_2008} indicates that these modifications result in 20\% reduction in computational cost 
without significant changes on physical performance.

The details of formulation are found in Section 8.3.2 of SCALE description \citep{scale_library}.


\subsection{Large-scale condensation scheme}
%
If the moisture content exceeds the saturation value, the excess moisture is converted to cloud water.
The latent heat is released to the atmosphere, and the cloud water is removed with a certain timescale.
The typical timescale is set to about O($10^3$) seconds when horizontal grid spacing is O($10^2$) km.

